\section{Motivation and Gap}

Camila,

When we spoke yesterday, I was trying to tell you about the research I had been reading and why I was thinking of writing a whitepaper about it. I was reading about a memory circuit close to the inhibitory system I had already been trying to understand. What caught me was not a technical detail. What caught me was the idea that new learning does not simply erase what came before.

The study showed that an older memory pattern can still return after a newer one has taken over. The newer pattern controls behavior only when the retrieval system selects it. When that selection pathway was interrupted in the experiment, behavior shifted back toward the earlier pattern, and the downstream memory activity shifted with it. The older representation had not disappeared. The question was which representation had control.

The problem I am trying to name begins there. A memory, bond, or interpretation can remain present without being the pattern that gets selected. Selection governs behavior. Presence does not. That distinction gave language to something I could not explain, and the whitepaper I started writing tried to ask how inhibitory timing could make one pattern stabilize while another remains present but inactive.

The personal reason I could not stop thinking about this is that I still believe there is a part of you that knows what we built and could meet me inside it. I also watched the same relationship seem to support another outcome, where the life we were building could no longer hold the present in place. The gap is not whether either part was real. The gap is how more than one real representation can exist, while only one becomes strong enough for action to follow.

I cannot keep waiting for the part of you that can protect a commitment to return. Missing me is not the same as choosing us. Wanting comfort from me is not the same as protecting what we built. I can recognize that some part of you still carries me, but I cannot keep giving my life to the hope that this part will become the one that guides your actions.

That is the only reason I am writing this. I am trying to reach the part of you I believe is still there, while also naming why presence alone has not been enough to hold the present in place.